\chapter{Espacios completos y completación}\label{cap:nuevo03}\label{chap:08}\label{chap:09}
Una sucesión de Cauchy contiene información interna: sus términos terminan siendo tan próximos como se quiera. La completitud asegura que esa información corresponde a un punto del propio espacio, y lo hace para todas las sucesiones de Cauchy. Comenzaremos por reconocer esta propiedad en espacios de sucesiones y en subespacios cerrados; después construiremos una completación para cualquier espacio métrico, conservando las distancias y haciendo que los puntos originales sean densos. La existencia de esa construcción y su unicidad relacionan los ejemplos concretos con una herramienta general. Usaremos la continuidad uniforme del capítulo~\ref{chap:07}, la propiedad arquimediana de $\R$ y las desigualdades elementales de las normas $\ell^p$.

\section{Límites que permanecen dentro del espacio}
\label{lectura:08}

\begin{proposicion}\label{ch08:cauchy-basico}
Toda sucesión convergente es de Cauchy. Toda sucesión de Cauchy es acotada. Si una sucesión de Cauchy tiene una subsucesión convergente a $x$, la sucesión completa converge a $x$.
\end{proposicion}
\begin{proof}
Si $x_n\to x$, dado $\eps>0$ elegimos $N$ tal que $d(x_n,x)<\eps/2$ para $n\geq N$. Entonces $d(x_n,x_m)<\eps$ para $n,m\geq N$. Si $(x_n)$ es de Cauchy, elegimos $N$ con $d(x_n,x_N)<1$ para $n\geq N$. El máximo de $1,d(x_1,x_N),\ldots,d(x_{N-1},x_N)$ es una cota de todas las distancias al centro $x_N$.

Por último, supongamos $x_{n_j}\to x$. Elegimos $N$ con $d(x_n,x_m)<\eps/2$ para $n,m\geq N$, y después $j$ con $n_j\geq N$ y $d(x_{n_j},x)<\eps/2$. Para cualquier $n\geq N$ se tiene $d(x_n,x)\leq d(x_n,x_{n_j})+d(x_{n_j},x)<\eps$.
\end{proof}

\begin{definicion}
Un espacio métrico es \emph{completo} si toda sucesión de Cauchy converge en él. Un espacio normado completo para la distancia $d(x,y)=\norm{x-y}$ se llama \emph{espacio de Banach}.
\end{definicion}
\index{completitud}\index{Banach, espacio de}

\begin{ejemplo}
En $\R\setminus\{0\}$, la sucesión $1/n$ es de Cauchy pero no converge en el espacio. De converger allí a $a$, también convergería a $a$ en $\R$, y la unicidad del límite obligaría a $a=0$. En cambio, todo espacio uniformemente discreto es completo: si los puntos distintos distan al menos $c>0$, la condición de Cauchy con tolerancia $c$ obliga a que la sucesión sea eventualmente constante. La mera topología discreta no garantiza completitud: $\{1/n:n\in\N\}$, con su distancia heredada, vuelve a ser un contraejemplo.
\end{ejemplo}

\section{Un ejemplo aritmético: la distancia p-ádica}
Fijemos un primo $p$. Todo racional no nulo se escribe como $r=p^k a/b$, donde $k\in\Z$, los enteros $a,b$ no son divisibles por $p$ y $b\ne0$. El exponente $k$ es único: es la potencia de $p$ del numerador menos la del denominador de cualquier fracción que represente a $r$. Los enteros $a,b$ no son únicos si no se exige además que la fracción esté reducida. Definimos $v_p(r)=k$, $|r|_p=p^{-k}$ y $|0|_p=0$. La distancia es
\[
 d_p(r,s)=|r-s|_p.
\]
Su propiedad decisiva, demostrada en el ejercicio~\ref{ej:lectura-8-1}, es
\[
 |r+s|_p\leq\max\{|r|_p,|s|_p\}.
\]
\index{metrica p-adica@métrica p-ádica}\index{ultrametrica@ultramétrica}
En esta geometría, ser divisible por una potencia grande de $p$ significa estar cerca de cero.

\begin{ejemplo}\label{ch08:padica-incompleta}
Para una sucesión de dígitos $a_j\in\{0,1,\ldots,p-1\}$ definamos $s_n=\sum_{j=0}^n a_jp^j$. Si $m>n$, la diferencia $s_m-s_n$ es divisible por $p^{n+1}$ y
\[
 |s_m-s_n|_p\leq p^{-(n+1)}\longrightarrow0.
\]
Se necesita una desigualdad, pues puede haber dígitos nulos adicionales o incluso diferencia cero. Por tanto $(s_n)$ es de Cauchy.

Hay incontables sucesiones de dígitos: si se enumeraran, podríamos cambiar el dígito de posición $n$ de la sucesión número $n$ y construir una que no aparece en la lista. Además, dos sucesiones distintas no pueden dar el mismo límite racional. Si su primer dígito distinto ocupa la posición $k$, para todo $n\geq k$ la diferencia entre sus sumas parciales tiene valoración exactamente $k$, pues el coeficiente de $p^k$ no es divisible por $p$. Las distancias entre esas sumas parciales son entonces constantemente $p^{-k}$; si ambas convergieran al mismo racional, esas distancias tenderían a cero. Como $\Q$ es numerable, no todas las sucesiones construidas pueden converger en $\Q$. Esto prueba que $(\Q,d_p)$ no es completo.

Algunas sí convergen: si $a_j=p-1$ para todo $j$, entonces $s_n=p^{n+1}-1$ y $|s_n-(-1)|_p=p^{-(n+1)}\to0$.
\end{ejemplo}

\section{Completitud de los espacios de sucesiones}
Recordemos que, para $1\leq p<\infty$,
\[
 \ell^p=\left\{x=(x_k):\sum_{k=1}^\infty|x_k|^p<\infty\right\},
 \qquad \norm{x}_p=\left(\sum_{k=1}^\infty|x_k|^p\right)^{1/p}.
\]
Las coordenadas pueden ser reales o complejas. En el segundo caso, la completitud de los escalares se obtiene aplicando la de $\R$ a partes reales e imaginarias.

\begin{teorema}\label{ch08:lp-completo}
Para cada $1\leq p<\infty$, el espacio $\ell^p$ es de Banach.
\end{teorema}
\begin{proof}
Sea $(x^{(n)})$ una sucesión de Cauchy en $\ell^p$, donde $x^{(n)}=(x^{(n)}_k)_{k\geq1}$. Distinguimos el índice $n$, que enumera vectores, del índice $k$, que enumera coordenadas. La estimación
\[
 |x^{(n)}_k-x^{(m)}_k|\leq\norm{x^{(n)}-x^{(m)}}_p
\]
muestra que cada sucesión escalar $(x^{(n)}_k)_n$ es de Cauchy. Definimos $x_k=\lim_nx^{(n)}_k$ y $x=(x_k)$.

Primero verificamos que $x$ pertenece a $\ell^p$. La sucesión de vectores es acotada, así que existe $M$ con $\norm{x^{(n)}}_p\leq M$ para todo $n$. Para cada entero $K$,
\[
 \sum_{k=1}^K|x_k|^p
 =\lim_{n\to\infty}\sum_{k=1}^K|x^{(n)}_k|^p\leq M^p.
\]
Aquí sólo se ha pasado el límite a través de una suma \emph{finita}. Las sumas parciales de la serie no negativa están acotadas y, por el axioma del supremo, convergen a un número no mayor que $M^p$.

Ahora probamos convergencia en norma. Dado $\eps>0$, existe $N$ tal que $\norm{x^{(n)}-x^{(m)}}_p<\eps/2$ siempre que $n,m\geq N$. Fijamos $n\geq N$ y $K\in\N$, y hacemos tender $m$ a infinito en
\[
 \sum_{k=1}^K|x^{(n)}_k-x^{(m)}_k|^p<(\eps/2)^p.
\]
Resulta $\sum_{k=1}^K|x^{(n)}_k-x_k|^p\leq(\eps/2)^p$. Tomando el supremo sobre $K$ obtenemos $\norm{x^{(n)}-x}_p\leq\eps/2<\eps$. El índice $N$ permanece fijo durante este argumento; no se hace tender a infinito un umbral que ya depende de $\eps$.
\end{proof}
La estrategia tiene dos partes que conviene distinguir: construir un candidato coordenada a coordenada y probar después una estimación uniforme en el número de coordenadas. La convergencia de cada coordenada, por sí sola, no implica convergencia en norma: los vectores canónicos $e_n$ convergen coordenada a coordenada a cero, pero $\norm{e_n}_p=1$.

\begin{teorema}[Subespacios completos]\label{ch08:cerrado-completo}
Sea $Y\subset X$ con la distancia heredada. Si $Y$ es completo, entonces es cerrado en $X$. Si $X$ es completo y $Y$ es cerrado en $X$, entonces $Y$ es completo.
\end{teorema}
\begin{proof}
Si $x\in\cl Y$, podemos elegir $y_n\in Y$ con $d(y_n,x)<1/n$. La sucesión es de Cauchy y, por completitud de $Y$, converge a algún $y\in Y$. También converge a $x$ en $X$, por lo que $x=y\in Y$. Esto prueba que $Y$ contiene su clausura. Si $Y$ es vacío, ya es cerrado.

Para la segunda afirmación, una sucesión de Cauchy en $Y$ lo es en $X$ porque las distancias son las mismas. Tiene, pues, un límite $x\in X$. Al ser $Y$ cerrado, todo límite de una sucesión suya pertenece a $Y$; así, $x\in Y$.
\end{proof}

\section{Ejercicios sobre completitud y espacios de sucesiones}
\begin{ejercicio}[Lectura 8, ejercicio 1]\label{ej:lectura-8-1}
Demuestre que $d_p$ es una métrica en $\Q$ y, más precisamente, una ultramétrica.
\end{ejercicio}
\begin{solucion}
La factorización prima garantiza la unicidad de $v_p(r)$ y $v_p(rs)=v_p(r)+v_p(s)$ para racionales no nulos. Sean $r,s\ne0$ y supongamos $v_p(r)=k\leq v_p(s)=l$. Escribiendo $r=p^ka/b$ y $s=p^lc/d$ con denominadores y numeradores no divisibles por $p$, obtenemos
\[
 r+s=p^k\frac{ad+p^{l-k}bc}{bd}.
\]
El denominador no es divisible por $p$, y el numerador es un entero. Si no es cero, su valoración es no negativa; de aquí $v_p(r+s)\geq k$. Si es cero, $|r+s|_p=0$. En ambos casos
\[
 |r+s|_p\leq p^{-k}=\max\{|r|_p,|s|_p\}.
\]
Si alguno de los sumandos es cero, la desigualdad se verifica directamente. También $|r|_p\geq0$, con igualdad exactamente para $r=0$, y $|-r|_p=|r|_p$. Por tanto $d_p$ es positiva, separa puntos y es simétrica. Aplicando la desigualdad anterior a $(r-t)+(t-s)$ obtenemos
\[
 d_p(r,s)\leq\max\{d_p(r,t),d_p(t,s)\}\leq d_p(r,t)+d_p(t,s).
\]
La primera desigualdad es la propiedad ultramétrica y la segunda da la desigualdad triangular usual.
\end{solucion}

\begin{ejercicio}[Lectura 8, ejercicio 2]\label{ej:lectura-8-2}
Demuestre la equivalencia entre el axioma del supremo y la completitud de $\R$ con su distancia usual.
\end{ejercicio}
\begin{solucion}
Trabajamos con las propiedades de cuerpo ordenado arquimediano de los reales; el contenido del ejercicio es comparar las dos formulaciones de su propiedad de completitud.

Supongamos el axioma del supremo y sea $(x_n)$ de Cauchy. Está acotada. Para cada $N$, definamos $a_N=\inf_{n\geq N}x_n$ y $b_N=\sup_{n\geq N}x_n$. El ínfimo existe aplicando el axioma del supremo al conjunto de los opuestos. La sucesión $(a_N)$ crece, $(b_N)$ decrece y $a_N\leq b_N$. Para cualquier $\eps>0$, la condición de Cauchy aplicada con $\eps/2$ implica que, a partir de cierto $N$, todos los términos de la cola están en un intervalo de longitud $\eps$; por ello $b_N-a_N\leq\eps$. Sea $x=\sup_Na_N$. Para todo $N$, se tiene $a_N\leq x\leq b_N$: la segunda desigualdad se debe a que $b_N$ es una cota superior de todos los $a_M$, también de los correspondientes a $M<N$. Si $n\geq N$, tanto $x_n$ como $x$ están en $[a_N,b_N]$, de modo que $|x_n-x|\leq b_N-a_N\to0$.

Recíprocamente, supongamos que toda sucesión real de Cauchy converge. Sea $A\ne\varnothing$ acotado superiormente. Si tiene máximo, ése es su supremo. En otro caso elegimos $a_0\in A$ y una cota superior $b_0>a_0$. La sucesión de intervalos se construye por bisección: si $m_n=(a_n+b_n)/2$ es cota superior de $A$, ponemos $(a_{n+1},b_{n+1})=(a_n,m_n)$; si no lo es, ponemos $(a_{n+1},b_{n+1})=(m_n,b_n)$. Aquí $a_n$ no tiene por qué pertenecer a $A$, pero nunca es cota superior, y $b_n$ siempre lo es. Los intervalos están anidados y su longitud es $(b_0-a_0)2^{-n}$. Esto hace que ambas sucesiones sean de Cauchy y que tengan el mismo límite $s$.

Todo $a\in A$ satisface $a\leq b_n$ para cada $n$, luego $a\leq s$. Si $t<s$, para algún $n$ se tiene $t<a_n$, y como $a_n$ no es cota superior existe $a\in A$ con $a>a_n>t$. Así, ningún número menor que $s$ es cota superior. Hemos probado $s=\sup A$. Si se desean puntos de $A$ como en la sugerencia original, elegimos $c_n\in A$ con $a_n<c_n\leq b_n$; entonces también $c_n\to s$.
\end{solucion}

\begin{ejercicio}[Lectura 8, ejercicio 3]\label{ej:lectura-8-3}
Demuestre que $\ell^\infty$ es completo.
\end{ejercicio}
\begin{solucion}
Sea $(x^{(n)})$ de Cauchy para $\norm{x}_\infty=\sup_k|x_k|$. Cada coordenada es de Cauchy, porque $|x^{(n)}_k-x^{(m)}_k|\leq\norm{x^{(n)}-x^{(m)}}_\infty$. Definamos $x_k=\lim_nx^{(n)}_k$. Existe $M$ con $\norm{x^{(n)}}_\infty\leq M$, y al pasar al límite para cada $k$ resulta $|x_k|\leq M$; así $x\in\ell^\infty$. Dado $\eps>0$, elegimos $N$ con $\norm{x^{(n)}-x^{(m)}}_\infty<\eps/2$ para $n,m\geq N$. Fijamos $n\geq N$ y dejamos $m\to\infty$: para cada $k$, $|x^{(n)}_k-x_k|\leq\eps/2$. Tomar el supremo en $k$ produce $\norm{x^{(n)}-x}_\infty\leq\eps/2<\eps$.
\end{solucion}

\begin{ejercicio}[Lectura 8, ejercicio 4]\label{ej:lectura-8-4}
Pruebe que los subespacios $c$ y $c_0$ de sucesiones convergentes y de sucesiones convergentes a cero son espacios de Banach con la norma uniforme.
\end{ejercicio}
\begin{solucion}
Ambos son subespacios lineales de $\ell^\infty$: las combinaciones lineales preservan la existencia del límite y, para $c_0$, preservan el límite cero. Por el teorema~\ref{ch08:cerrado-completo}, basta probar que son cerrados.

Sea $x^{(n)}\in c$ y $\norm{x^{(n)}-x}_\infty\to0$. Para $\eps>0$ elegimos $n$ con $\norm{x^{(n)}-x}_\infty<\eps/3$. Como la sucesión escalar $k\mapsto x^{(n)}_k$ converge, existe $K$ con $|x^{(n)}_k-x^{(n)}_j|<\eps/3$ para $k,j\geq K$. Por tanto
\[
 |x_k-x_j|\leq|x_k-x^{(n)}_k|+|x^{(n)}_k-x^{(n)}_j|+|x^{(n)}_j-x_j|<\eps.
\]
La sucesión de coordenadas de $x$ es de Cauchy en los escalares y converge; de ahí $x\in c$.

Si cada $x^{(n)}\in c_0$, elegimos $n$ con $\norm{x^{(n)}-x}_\infty<\eps/2$ y luego $K$ con $|x^{(n)}_k|<\eps/2$ para $k\geq K$. Entonces $|x_k|<\eps$ para $k\geq K$, luego $x\in c_0$.
\end{solucion}

\begin{ejercicio}[Lectura 8, ejercicio 5]\label{ej:lectura-8-5}
Para espacios métricos no vacíos, pruebe que el producto de dos espacios es completo si y sólo si cada factor lo es. Extienda el resultado a todo producto finito.
\end{ejercicio}
\begin{solucion}
Usamos la métrica máxima $D((x,y),(x',y'))=\max\{d(x,x'),\rho(y,y')\}$. Una sucesión es de Cauchy para $D$ exactamente cuando sus dos coordenadas lo son: una cota para $D$ acota ambas distancias, y dos umbrales para las coordenadas se reúnen tomando su máximo. Si ambos factores son completos, las coordenadas convergen a $x,y$, y el máximo de sus distancias a esos límites tiende a cero.

Si el producto es completo, fijamos $y_0\in Y$ y consideramos una sucesión de Cauchy $(x_n)$ de $X$. Los pares $(x_n,y_0)$ forman una sucesión de Cauchy que converge a algún $(x,y)$, y la primera coordenada da $x_n\to x$. Fijando $x_0$ se demuestra lo mismo para $Y$. Para $n$ factores, la métrica $\max_{1\leq j\leq n}d_j(x_j,y_j)$ permite repetir exactamente las dos implicaciones: tomamos el máximo de los $n$ umbrales y, para la necesidad, fijamos las otras $n-1$ coordenadas. La métrica suma y la euclídea de las distancias son equivalentes mediante cotas constantes a la métrica máxima; por ello tienen las mismas sucesiones de Cauchy. La no vacuidad evita que un producto vacío oculte un factor incompleto.
\end{solucion}

\begin{ejercicio}[Lectura 8, ejercicio 6]\label{ej:lectura-8-6}
Demuestre que el conjunto de Cantor, con la distancia heredada de $\R$, es completo y no tiene puntos aislados.
\end{ejercicio}
\begin{solucion}
Sea $C_0=[0,1]$ y sea $C_n$ el conjunto que queda después de eliminar los tercios medios abiertos durante las primeras $n$ etapas. Cada $C_n$ es una unión finita de intervalos cerrados de longitud $3^{-n}$, y $C=\bigcap_{n\geq0}C_n$ es cerrado en $\R$. Como $\R$ es completo, también lo es $C$.

Fijemos $x\in C$ y $\eps>0$. Elegimos $n$ con $3^{-n}<\eps$ y un intervalo componente $[a,b]$ de $C_n$ que contenga $x$. Ambos extremos pertenecen a todas las etapas posteriores: nunca se elimina un extremo al retirar un tercio medio abierto. Por tanto $a,b\in C$. Como $a\ne b$, al menos uno de ellos es distinto de $x$ y dista de $x$ a lo sumo $3^{-n}<\eps$. Ningún punto de $C$ es aislado.
\end{solucion}

\begin{ejercicio}[Lectura 8, ejercicio 7]\label{ej:lectura-8-7}
Demuestre que $C$ es incontable y que
\[
 C=\left\{\sum_{n=1}^\infty\frac{a_n}{3^n}:a_n\in\{0,2\}\right\}.
\]
\end{ejercicio}
\begin{solucion}
Para cada palabra $(a_1,\ldots,a_n)$ de ceros y doses, una componente de $C_n$ es
\[
 I(a_1,\ldots,a_n)=\left[\sum_{j=1}^n\frac{a_j}{3^j},\,
 \sum_{j=1}^n\frac{a_j}{3^j}+3^{-n}\right].
\]
Esto se verifica por inducción: escoger el tercio izquierdo agrega el dígito $0$ y escoger el derecho agrega el dígito $2$. Una serie con esos dígitos pertenece al intervalo indicado, pues su cola está entre cero y $\sum_{j>n}2/3^j=3^{-n}$. Pertenece entonces a todo $C_n$ y a $C$.

Recíprocamente, si $x\in C$, en cada etapa está en una componente y la componente de la etapa siguiente está contenida en la anterior. Sus elecciones izquierda/derecha proporcionan dígitos $a_n\in\{0,2\}$. La diferencia entre $x$ y la suma de sus primeros $n$ dígitos está entre cero y $3^{-n}$, por lo que $x$ es la serie correspondiente. No es necesario usar la expansión ternaria convencional en los extremos: por ejemplo, $1/3=0.02222\ldots$ usa sólo dígitos permitidos.

La correspondencia es inyectiva. Si dos sucesiones difieren por primera vez en $k$, la diferencia aportada allí tiene valor absoluto $2/3^k$, mientras que toda la cola posterior puede compensar a lo sumo $\sum_{j>k}2/3^j=1/3^k$. Las sumas difieren al menos en $1/3^k>0$. Finalmente, el conjunto de sucesiones de ceros y doses es incontable por el argumento diagonal: frente a una supuesta enumeración, cambiamos el dígito $n$ de la fila $n$. Su imagen inyectiva $C$ es, pues, incontable.
\end{solucion}

\section{Agregar límites sin alterar las distancias}
\label{lectura:09}
Los espacios $\ell^p$, $c$ y $c_0$ ofrecen modelos completos, mientras que los racionales p-ádicos muestran que una noción distinta de distancia puede exigir nuevos puntos. Buscamos ahora una construcción que agregue sistemáticamente los límites que faltan. Las distancias entre los puntos originales deben conservarse, y los puntos nuevos deben surgir como límites de puntos antiguos. La completación que construiremos satisface ambas exigencias y es única salvo una identificación que conserva distancias.

\begin{definicion}
Una aplicación $i:(X,d)\to(Y,\rho)$ es una \emph{isometría} si $\rho(i(x),i(y))=d(x,y)$ para todos $x,y\in X$. Una \emph{completación} de $X$ es un espacio completo $\widehat X$ junto con una isometría $i:X\to\widehat X$ cuya imagen es densa.
\end{definicion}
\index{isometria@isometría}\index{completacion@completación}
Una isometría es inyectiva: imágenes iguales dan distancia cero entre los puntos originales. Es $1$-Lipschitz y, por tanto, uniformemente continua; su inversa sobre la imagen también es una isometría. La inclusión de un subespacio con la distancia heredada es un ejemplo. Si $X$ ya es completo, $(X,\id_X)$ es una completación suya.

\section{Construcción mediante sucesiones de Cauchy}
\begin{teorema}\label{ch09:existencia}
Todo espacio métrico tiene una completación.
\end{teorema}
\begin{proof}
El espacio vacío es completo y se completa con él mismo. Supongamos $X\ne\varnothing$ y sea $\mathcal C(X)$ el conjunto de todas las sucesiones de Cauchy de $X$. Este conjunto no es vacío porque contiene todas las sucesiones constantes.

\textbf{Paso 1: una distancia provisional.} Para $u=(u_n)$ y $v=(v_n)$ ponemos
\[
 D(u,v)=\lim_{n\to\infty}d(u_n,v_n).
\]
Debemos demostrar que el límite existe. La desigualdad triangular, aplicada en los dos sentidos, produce
\[
 |d(u_n,v_n)-d(u_m,v_m)|\leq d(u_n,u_m)+d(v_n,v_m).
\]
Las dos sucesiones del lado derecho son de Cauchy; así, $(d(u_n,v_n))$ es una sucesión real de Cauchy. Su límite existe por la completitud de $\R$. Además, $D\geq0$, $D(u,u)=0$ y $D(u,v)=D(v,u)$. Si $w=(w_n)$, de $d(u_n,v_n)\leq d(u_n,w_n)+d(w_n,v_n)$ y la preservación de desigualdades al pasar al límite resulta $D(u,v)\leq D(u,w)+D(w,v)$. Por ello $D$ es una pseudométrica; puede anularse sobre sucesiones diferentes, por ejemplo sobre dos sucesiones que converjan al mismo punto.

\textbf{Paso 2: identificar las representaciones del mismo punto.} Definimos $u\sim v$ cuando $D(u,v)=0$. La reflexividad y la simetría provienen de las de $D$; la transitividad se sigue de $D(u,w)\leq D(u,v)+D(v,w)$. Sea
\[
 \widehat X=\mathcal C(X)/\!\sim,
 \qquad \widehat d([u],[v])=D(u,v).
\]
Esta fórmula no depende de representantes: si $u\sim u'$ y $v\sim v'$, entonces
\[
 |D(u,v)-D(u',v')|\leq D(u,u')+D(v,v')=0.
\]
Las propiedades de pseudométrica descienden al cociente y, además, $\widehat d([u],[v])=0$ implica $[u]=[v]$. Tenemos una métrica verdadera.

\textbf{Paso 3: insertar $X$ y probar densidad.} Definimos $i(x)$ como la clase de la sucesión constante $(x,x,\ldots)$. Entonces $\widehat d(i(x),i(y))=\lim_nd(x,y)=d(x,y)$, por lo que $i$ es una isometría. Si $u=[(u_n)]$, se cumple
\[
 \widehat d(i(u_n),u)=\lim_{k\to\infty}d(u_n,u_k)\longrightarrow0
 \quad(n\to\infty).
\]
En efecto, dado $\eps>0$, la condición de Cauchy permite que $d(u_n,u_k)<\eps/2$ para $n,k\geq N$; fijando $n\geq N$ y pasando $k$ al límite, queda una distancia no mayor que $\eps/2<\eps$. Esto prueba que $i(X)$ es denso.

\textbf{Paso 4: completar la construcción.} Sea $(U_n)$ una sucesión de Cauchy en $\widehat X$. Por la densidad ya probada, para cada $n$ elegimos $z_n\in X$ con
\[
 \widehat d(i(z_n),U_n)<2^{-n}.
\]
Estas elecciones se hacen una sola vez, con tolerancias predeterminadas. No dependen del $\eps$ que aparecerá después al comprobar la convergencia. Para $n,m\in\N$,
\[
 d(z_n,z_m)\leq 2^{-n}+\widehat d(U_n,U_m)+2^{-m}.
\]
Dado $\eps>0$, elegimos $N$ de modo que los dos términos geométricos sumen menos que $\eps/2$ para $n,m\geq N$ y que $\widehat d(U_n,U_m)<\eps/2$. Así $(z_n)$ es de Cauchy en $X$. Sea $Z=[(z_n)]\in\widehat X$. El paso 3 aplicado a esta sucesión da $i(z_n)\to Z$. Finalmente,
\[
 \widehat d(U_n,Z)\leq\widehat d(U_n,i(z_n))+\widehat d(i(z_n),Z)
 <2^{-n}+\widehat d(i(z_n),Z)\longrightarrow0.
\]
Toda sucesión de Cauchy de $\widehat X$ converge allí; por tanto es completo.
\end{proof}

\begin{observacion}[La elección diagonal y sus cuantificadores]
Si cada $U_n$ se representa por una sucesión $(u_{n,k})_k$, el punto $z_n$ del paso 4 puede tomarse como $u_{n,k_n}$ con $k_n$ suficientemente grande. Se puede imponer además $k_n>k_{n-1}$, porque la aproximación vale para todos los índices de una cola. Lo esencial es escoger $\widehat d(i(u_{n,k_n}),U_n)<2^{-n}$, y no un error fijo $\eps$ antes de definir la diagonal: una sucesión candidata al límite debe estar definida independientemente de la tolerancia con la que más tarde se compruebe que converge.
\end{observacion}

\section{Dos modelos y una precaución sobre las normas}
\begin{ejemplo}[Aproximación de funciones discontinuas]\leavevmode\par
En $C([0,1])$, la fórmula $\norm{f}_1=\int_0^1|f(t)|\,dt$ define una norma. Si una función continua no es cero, su valor absoluto es mayor que una constante positiva en algún intervalo de longitud positiva, y su integral es positiva. El valor absoluto es indispensable: $\int f$ por sí sola no es una norma.

Este espacio normado no es completo. Para $n\geq3$, definamos $f_n$ igual a cero en $[0,1/2]$, lineal de cero a uno en $[1/2,1/2+1/n]$ e igual a uno en el intervalo restante. Si $h$ es cero en $[0,1/2]$ y uno en $(1/2,1]$, entonces
\[
 \int_0^1|f_n-h|\leq1/n,
 \quad\text{y por tanto}\quad
 \norm{f_n-f_m}_1\leq1/n+1/m.
\]
Si existiera un límite continuo $f$ en esta norma, la misma desigualdad triangular daría $\int|f-h|=0$. En cualquier intervalo contenido en $(0,1/2)$, la continuidad obligaría a $f=0$: un valor no nulo produciría una integral positiva en una vecindad. De igual forma, $f=1$ en $(1/2,1)$. Esto contradice la continuidad en $1/2$.

El espacio de clases de funciones Lebesgue-integrables $L^1([0,1])$ proporcionará una realización de esta completación en el volumen de medida e integración. Aquí no se presupone esa teoría: la completación por sucesiones ya está construida. Para un compacto arbitrario $K\subset\R^n$, la integral puede ser sólo una seminorma sobre $C(K)$; por ejemplo, si $K=\{0\}$, todas las integrales respecto del volumen son cero. Por ello la afirmación sobre $C(K)$ requiere hipótesis sobre la medida y su soporte, o pasar primero al cociente por las funciones de seminorma cero.
\end{ejemplo}

\begin{ejemplo}[Los números p-ádicos]\label{ch09:qp}
La completación de $(\Q,d_p)$ se denota por $\Q_p$. Para una clase $r=[(r_n)]$, definimos $|r|_p=\widehat d(r,0)=\lim_n|r_n|_p$. La desigualdad ultramétrica pasa al límite, de modo que sigue siendo válida en la completación. Además de ser un espacio completo, conserva las operaciones de cuerpo. Veamos por qué. Si $(r_n)$ y $(s_n)$ son de Cauchy, sus valores absolutos p-ádicos están acotados, y
\[
 |r_ns_n-r_ms_m|_p\leq
 \max\{|r_n|_p|s_n-s_m|_p,\,|s_m|_p|r_n-r_m|_p\}.
\]
Esto prueba que los productos forman una sucesión de Cauchy; para la suma, $|(r_n+s_n)-(r_m+s_m)|_p$ está acotado por el máximo de las dos diferencias correspondientes. Las mismas estimaciones muestran que cambiar representantes por sucesiones equivalentes no cambia la clase de su suma o producto. Las identidades de cuerpo se conservan porque valen término a término. También $|rs|_p=|r|_p|s|_p$, pues la identidad racional $|r_ns_n|_p=|r_n|_p|s_n|_p$ pasa al límite de números reales.

Falta justificar la inversa. Para una clase $r\ne0$, su distancia a cero es un número $c>0$. Si $(r_n)$ la representa, entonces $|r_n|_p\to c$, de modo que $|r_n|_p\geq c/2$ en una cola. Allí
\[
 |r_n^{-1}-r_m^{-1}|_p
 =\frac{|r_n-r_m|_p}{|r_n|_p|r_m|_p}
 \leq \frac{4}{c^2}|r_n-r_m|_p.
\]
La sucesión de inversas es de Cauchy y su clase multiplica a $r$ para dar $1$. Los términos iniciales se pueden reemplazar sin alterar la clase. Así $\Q_p$ es un cuerpo.

Todo elemento tiene una expansión
\[
 r=\sum_{k\geq k_0}a_kp^k,
 \qquad k_0\in\Z,\quad a_k\in\{0,\ldots,p-1\},
\]
entendida como límite p-ádico de sus sumas parciales. Damos la razón de la existencia. Primero supongamos $|r|_p\leq1$ y tomemos racionales $r_n\to r$. Para cada $m\geq1$, elegimos uno con $|r-r_n|_p<p^{-m}$ y $|r_n|_p\leq1$. Podemos escribirlo $r_n=a/b$ con $p\nmid b$. Como $b$ tiene inversa módulo $p^m$, existe un único entero $c_m\in\{0,\ldots,p^m-1\}$ con $bc_m\equiv a\pmod{p^m}$. Por tanto $|r-c_m|_p\leq p^{-m}$. La unicidad de estos residuos, y la desigualdad ultramétrica, implican $c_{m+1}\equiv c_m\pmod{p^m}$; en consecuencia $c_{m+1}=c_m+a_mp^m$ para un dígito $a_m$. Partiendo de $c_1=a_0$, obtenemos la expansión y la cota $|r-c_m|_p\leq p^{-m}\to0$.

Para un $r$ arbitrario, su norma es finita y podemos multiplicarlo por una potencia positiva de $p$ hasta que su norma sea a lo sumo $1$; deshacer esa multiplicación produce el índice inicial $k_0$ posiblemente negativo. La unicidad de los dígitos se demuestra como en el ejemplo~\ref{ch08:padica-incompleta}: el primer dígito distinto determina la valoración de la diferencia y no puede ser cancelado por términos de mayor potencia. Este argumento muestra también cómo la completación agrega puntos que la aritmética racional no contenía.
\end{ejemplo}

\section{Unicidad y extensión: ejercicio resuelto}
\begin{ejercicio}[Lectura 9, ejercicio 1]\label{ej:lectura-9-1}\label{ch09:unicidad}
Sean $(Y,i)$ y $(Y',i')$ dos completaciones de $X$. Demuestre que existe una isometría sobreyectiva $T:Y\to Y'$ tal que $T\circ i=i'$.
\end{ejercicio}
\begin{solucion}
Si $X=\varnothing$, la densidad fuerza $Y=Y'=\varnothing$, y la afirmación es inmediata. En otro caso, para $y\in Y$ elegimos $x_n\in X$ con $i(x_n)\to y$. Como $i$ es isometría, $(x_n)$ es de Cauchy en $X$, y $(i'(x_n))$ lo es en el espacio completo $Y'$. Definimos
\[
 T(y)=\lim_{n\to\infty}i'(x_n).
\]
Si otra sucesión $z_n$ satisface $i(z_n)\to y$, entonces
\[
 d(x_n,z_n)=\rho(i(x_n),i(z_n))
 \leq\rho(i(x_n),y)+\rho(y,i(z_n))\longrightarrow0.
\]
La distancia entre $i'(x_n)$ e $i'(z_n)$ es la misma; sus límites coinciden y $T$ está bien definida. Para $y,z\in Y$, elegimos aproximaciones $i(x_n)\to y$, $i(z_n)\to z$ y obtenemos
\[
 \rho'(T(y),T(z))=\lim_nd(x_n,z_n)=\rho(y,z).
\]
Aquí la continuidad de la distancia se justifica por $|\rho(a,b)-\rho(c,d)|\leq\rho(a,c)+\rho(b,d)$. Por tanto $T$ es una isometría. La sucesión constante $x_n=x$ muestra $T(i(x))=i'(x)$.

Para probar sobreyectividad, sea $y'\in Y'$. Elegimos $x_n\in X$ con $i'(x_n)\to y'$. Entonces $(i(x_n))$ es de Cauchy y converge en $Y$ a cierto $y$. Por la definición anterior, $T(y)=y'$. Finalmente, la isometría que cumple la identidad requerida es única: es continua y queda determinada en la imagen densa $i(X)$; si $S$ es otra, para $i(x_n)\to y$ tenemos $S(y)=\lim_ni'(x_n)=T(y)$.
\end{solucion}

\begin{corolario}[Propiedad de extensión]\label{ch09:extension}
Toda aplicación uniformemente continua $f:X\to Z$, con $Z$ completo, se extiende de manera única a una aplicación uniformemente continua $\widehat f:\widehat X\to Z$ tal que $\widehat f\circ i=f$.
\end{corolario}
\begin{proof}
La aplicación $f\circ i^{-1}:i(X)\to Z$ es uniformemente continua porque $i^{-1}$ conserva distancias. Aplicamos la construcción y las estimaciones del ejercicio~\ref{ch07:ej-extension} al conjunto denso $i(X)\subset\widehat X$. Su clausura es todo $\widehat X$, por lo que obtenemos exactamente la extensión solicitada y su unicidad.
\end{proof}

\section{Completitud, construcción y unicidad}
Un punto nuevo es una clase de aproximaciones que terminan siendo indistinguibles en distancia. La densidad evita añadir puntos ajenos a esa idea, y la unicidad prueba que el resultado no depende de la realización concreta escogida. La completitud convierte después construcciones infinitas en puntos efectivos del espacio.
